\chapter*{ملخص}
يقدم هذا المشروع تصميماً وتنفيذاً متكاملاً لمنظومة ذراع روبوتية ذات خمس درجات حرية (\textenglish{5-DOF}) مخصصة لمهام التقاط ونقل الأشياء عبر الأوامر الصوتية، مما يمثل حلاً هندسياً متطوراً في مجال التفاعل البشري-الروبوتي (\textenglish{Human-Robot Interaction - HRI}). يرتكز العمل على بنية تحكم هرمية تبدأ من النمذجة الرياضية الحركية والديناميكية، وتمر عبر تصميم خوارزميات تحكم متينة، وتستقر عند طبقة الإشراف الصوتي الذاتي.

في الجانب الحركي، تم الاعتماد على اتفاقية دينافيت-هارتنبرغ (\textenglish{DH}) لبناء نموذج الحركيات الأمامية، بينما تم اشتقاق معادلات الحركيات العكسية تحليلياً بفصل الموضع عن التوجيه عبر نقطة مركز المعصم. وفي الجانب الديناميكي، تم صياغة معادلات الحركة باستخدام صياغة لاغرانج (\textenglish{Lagrange}) لاستخراج مصفوفات العطالة ومتجهات الجاذبية، مع إهمال قوى كوريوليس والقوى الطاردة المركزية نظراً لطبيعة السرعات المنخفضة، وجرى تصدير النموذج الميكانيكي من برنامج \textenglish{SolidWorks} إلى بيئة \textenglish{Simscape Multibody} في \textenglish{MATLAB/Simulink} لإجراء المحاكاة الرقمية والتحقق من الاستجابة الديناميكية للنظام تحت تأثير الأحمال والقصور الذاتي المتغير.

يتضمن شق التحكم تصميم ومقارنة ثلاث خوارزميات للتحكم منخفض المستوى (\textenglish{LLC}): متحكم \textenglish{PID} التقليدي المدعوم بتعويض الجاذبية، التحكم التكيفي المعتمد على العزم المحسوب (\textenglish{ACTC}) للتعامل مع تغيرات الحمولة المفاجئة، والتحكم بالوضع الانزلاقي (\textenglish{SMC}) لضمان المتانة العالية ضد الاضطرابات مع استخدام دالة التشبع للحد من ظاهرة التشذر (\textenglish{Chattering}) الضارة بالمشغلات، متبوعاً بإغلاق حلقة التحكم الكارتزي عالي المستوى (\textenglish{HLC}) لربط مراجع الموضع الكارتزية بالمراجع المفصلية.

وأخيراً، تم تتويج المنظومة بدمج نظام تفاعل صوتي مدمج ومستقل يعتمد على حزمة \textenglish{Vosk} للتعرف على الكلام دون اتصال بالإنترنت (\textenglish{Offline Processing})، حيث تقوم الخوارزمية بتحويل الأمر الصوتي إلى نص وتصفيته لاستخلاص الكلمات المفتاحية وترجمتها لحظياً إلى إحداثيات مرجعية كارتزية ($X_{\text{ref}}$) لتحديث مسار الذراع، مما يثبت هندسياً ورياضياً استقرار المنظومة الروبوتية المتكاملة وجاهزيتها الكاملة للتنفيذ والتحقق العملي بكفاءة عالية وبنية متينة.